\FigureLayoutDeclare{img/2003/beijing_liberal_autumn/q18_fig1.png}{11.0cm}{19bp 87bp 101bp 67bp}

\chapter{2003年北京卷（秋文）}
\section{单选题}
\begin{problem}
设集合 \(A = \{x \mid x^2 - 1 > 0\}\)，\(B = \{x \mid \log_2 x > 0\}\)，\(A \cap B\) 等于（\quad）

\choices
    {\(\{x \mid x > 1\}\)}
    {\(\{x \mid x > 0\}\)}
    {\(\{x \mid x < -1\}\)}
    {\(\{x \mid x < -1 \text{\ 或\ } x > 1\}\)}
\end{problem}

\begin{answer}
A
\end{answer}

\begin{solution}
由 \(x^2-1>0\) 得 \(x<-1\) 或 \(x>1\)，而 \(\log_2x>0\) 等价于 \(x>1\)，所以 \(A\cap B=\{x\mid x>1\}\)，故选 A.
\end{solution}

\begin{problem}
设 \(y_1 = 4^{0.9}\)，\(y_2 = 8^{0.44}\)，\(y_3 = \left(\frac{1}{2}\right)^{-1.5}\)，则（\quad）

\choices
    {\(y_3 > y_1 > y_2\)}
    {\(y_2 > y_1 > y_3\)}
    {\(y_1 > y_2 > y_3\)}
    {\(y_1 > y_3 > y_2\)}
\end{problem}

\begin{answer}
D
\end{answer}

\begin{solution}
将三数都化为以 \(2\) 为底的幂，有 \(y_1=4^{0.9}=2^{1.8}\)，\(y_2=8^{0.44}=2^{1.32}\)，\(y_3=(1/2)^{-1.5}=2^{1.5}\). 因为底数 \(2>1\)，指数大小规律与数值大小规律相同，且 \(1.8>1.5>1.32\)，所以 \(y_1>y_3>y_2\)，故选 D.
\end{solution}

\begin{problem}
“\(\cos 2\alpha = -\frac{\sqrt{3}}{2}\)”是“\(\alpha = k\pi + \frac{5\pi}{12}\)，（\(k \in \Z\)）”的（\quad）

\choices
    {必要非充分条件}
    {充分非必要条件}
    {充分必要条件}
    {既非充分又非必要条件}
\end{problem}

\begin{answer}
A
\end{answer}

\begin{solution}
由 \(\cos 2\alpha=-\dfrac{\sqrt{3}}{2}\)，得
\(2\alpha=\dfrac{5\pi}{6}+2k\pi\) 或 \(2\alpha=\dfrac{7\pi}{6}+2k\pi\)，其中 \(k\in\Z\). 因而
\(\alpha=\dfrac{5\pi}{12}+k\pi\) 或 \(\alpha=\dfrac{7\pi}{12}+k\pi\). 所以条件 \(\alpha=k\pi+\dfrac{5\pi}{12}\) 能推出 \(\cos 2\alpha=-\dfrac{\sqrt{3}}{2}\)，但后一个等式还允许 \(\alpha=k\pi+\dfrac{7\pi}{12}\)，不一定推出 \(\alpha=k\pi+\dfrac{5\pi}{12}\). 因此前一个等式是后一个条件的必要非充分条件，选 A.
\end{solution}

\begin{problem}
已知 \(\alpha\)，\(\beta\) 是平面，\(m\)，\(n\) 是直线. 下列命题中不正确的是（\quad）

\choices
    {若 \(m \parallel n\)，\(m \perp \alpha\)，则 \(n \perp \alpha\)}
    {若 \(m \parallel \alpha\)，\(\alpha \cap \beta = n\)，则 \(m \parallel n\)}
    {若 \(m \perp \alpha\)，\(m \perp \beta\)，则 \(\alpha \parallel \beta\)}
    {若 \(m \perp \alpha\)，\(m \subset \beta\)，则 \(\alpha \perp \beta\)}
\end{problem}

\begin{answer}
B
\end{answer}

\begin{solution}
命题 A 正确：若 \(m\parallel n\) 且 \(m\perp\alpha\)，则与 \(m\) 平行的直线 \(n\) 也垂直于平面 \(\alpha\). 命题 C 正确，因为垂直于同一直线的两个平面平行；命题 D 正确，因为平面 \(\beta\) 含有直线 \(m\)，而 \(m\perp\alpha\)，所以 \(\alpha\perp\beta\). 命题 B 不正确，例如取 \(\alpha:z=0\)，\(\beta:x=0\)，则 \(n=\alpha\cap\beta\) 是 \(y\) 轴；取 \(m:y=0\)，\(z=1\)，有 \(m\parallel\alpha\)，但 \(m\) 与 \(n\) 既不平行也不相交，故不能推出 \(m\parallel n\)，选 B.
\end{solution}

\begin{problem}
如图，直线 \(l: x - 2y + 2 = 0\) 过椭圆的左焦点 \(F_1\) 和一个顶点 \(B\)，该椭圆的离心率是（\quad）

\bitmapfigure[max width=0.42\linewidth]{img/2003/beijing_liberal_autumn/q5_fig1.png}

\choices
    {\(\frac{1}{5}\)}
    {\(\frac{2}{5}\)}
    {\(\frac{\sqrt{5}}{5}\)}
    {\(\frac{2\sqrt{5}}{5}\)}
\end{problem}

\begin{answer}
D
\end{answer}

\begin{solution}
由直线 \(l:x-2y+2=0\) 与 \(x\) 轴交于 \(F_1\)，令 \(y=0\)，得 \(F_1=(-2,0)\)，所以椭圆的左焦点参数 \(c=2\). 直线与 \(y\) 轴交于顶点 \(B\)，令 \(x=0\)，得 \(B=(0,1)\)，所以短半轴 \(b=1\). 对椭圆有 \(a^2=b^2+c^2=1+4=5\)，故离心率 \(e=c/a=2/\sqrt5=2\sqrt5/5\)，选 D.
\end{solution}

\begin{problem}
若 \(z \in \mathbf{C}\) 且 \(\left|z + 2 - 2\i\right| = 1\)，则 \(\left|z - 2 - 2\i\right|\) 的最小值是（\quad）

\choices
    {\(2\)}
    {\(3\)}
    {\(4\)}
    {\(5\)}
\end{problem}

\begin{answer}
B
\end{answer}

\begin{solution}
令 \(w=z+2-2\i\)，则 \(\lvert w\rvert=1\)，且
\(z-2-2\i=w-4\). 由三角不等式的逆不等式，
\(\lvert z-2-2\i\rvert=\lvert w-4\rvert\geqslant\bigl|\lvert4\rvert-\lvert w\rvert\bigr|=3\). 当 \(w=1\) 时，\(\lvert w\rvert=1\) 且 \(\lvert w-4\rvert=3\)，等号可以取得，所以最小值为 \(3\)，选 B.
\end{solution}

\begin{problem}
如果圆台的母线与底面成 \(60^{\circ}\) 角，那么这个圆台的侧面积与轴截面面积的比为（\quad）

\choices
    {\(2\pi\)}
    {\(\frac{3}{2}\pi\)}
    {\(\frac{2\sqrt{3}}{3}\pi\)}
    {\(\frac{1}{2}\pi\)}
\end{problem}

\begin{answer}
C
\end{answer}

\begin{solution}
设圆台上、下底面半径分别为 \(r\)、\(R\)，母线长为 \(l\)，高为 \(h\). 母线与底面成 \(60^\circ\) 角，所以 \(h=l\sin60^\circ=\sqrt3l/2\). 圆台侧面积为 \(S_1=\pi(R+r)l\)，轴截面是上、下底为 \(2r\)，\(2R\)、高为 \(h\) 的梯形，面积为 \(S_2=(R+r)h\). 因而 \(S_1/S_2=\pi l/h=2\pi/\sqrt3=2\sqrt3\pi/3\)，选 C.
\end{solution}

\begin{problem}
若数列 \(\{a_n\}\) 的通项公式是 \(a_n = \frac{3^{-n} + (-1)^n 3^{-n}}{2}\)，\(n = 1\)，\(2\)，\(\ldots\)，则 \(\displaystyle\lim_{n \to \infty}(a_1 + a_2 + \cdots + a_n)\) 等于（\quad）

\choices
    {\(\frac{1}{24}\)}
    {\(\frac{1}{8}\)}
    {\(\frac{1}{6}\)}
    {\(\frac{1}{2}\)}
\end{problem}

\begin{answer}
B
\end{answer}

\begin{solution}
当 \(n\) 为奇数时，\(1+(-1)^n=0\)，所以 \(a_n=0\)；当 \(n=2k\) 时，
\(a_{2k}=\dfrac{3^{-2k}+3^{-2k}}{2}=3^{-2k}=\left(\dfrac19\right)^k\). 因而部分和的极限为
\(\displaystyle\sum_{k=1}^{\infty}\left(\dfrac19\right)^k=\dfrac{\frac19}{1-\frac19}=\dfrac18\)，故选 B.
\end{solution}

\begin{problem}
从黄瓜，白菜，油菜，扁豆 \(4\) 种蔬菜品种中选出 \(3\) 种，分别种在不同土质的三块土地上，其中黄瓜必须种植，不同的种植方法共有（\quad）

\choices
    {\(24\) 种}
    {\(18\) 种}
    {\(12\) 种}
    {\(6\) 种}
\end{problem}

\begin{answer}
B
\end{answer}

\begin{solution}
黄瓜必须选出，再从白菜、油菜、扁豆中选出 \(2\) 种，有 \(\mathrm C_3^2=3\) 种选法. 选定的 \(3\) 种蔬菜分别种在三块不同土质的土地上，有 \(\mathrm A_3^3=3!=6\) 种安排方法. 所以不同的种植方法共有 \(\mathrm C_3^2\mathrm A_3^3=3\times6=18\) 种，选 B.
\end{solution}

\begin{problem}
某班试用电子投票系统选举班干部候选人. 全班 \(k\) 名同学都有选举权和被选举权，他们的编号分别为 \(1\)，\(2\)，\(\ldots\)，\(k\)，规定：同意按“1”，不同意（含弃权）按“0”，令
\(a_{ij} = \begin{cases}
1,& \text{第 } i\text{ 号同学同意第 }j\text{ 号同学当选},\\
0,& \text{第 } i\text{ 号同学不同意第 }j\text{ 号同学当选}.
\end{cases}\)
其中 \(i = 1\)，\(2\)，\(\ldots\)，\(k\)，且 \(j = 1\)，\(2\)，\(\ldots\)，\(k\)，则同时同意第 \(1\)，\(2\) 号同学当选的人数为（\quad）

\choices
    {\(a_{11} + a_{12} + \cdots + a_{1k} + a_{21} + a_{22} + \cdots + a_{2k}\)}
    {\(a_{11} + a_{21} + \cdots + a_{k1} + a_{12} + a_{22} + \cdots + a_{k2}\)}
    {\(a_{11}a_{12} + a_{21}a_{22} + \cdots + a_{k1}a_{k2}\)}
    {\(a_{11}a_{21} + a_{12}a_{22} + \cdots + a_{1k}a_{2k}\)}
\end{problem}

\begin{answer}
C
\end{answer}

\begin{solution}
对第 \(i\) 号同学，同时同意第 \(1\)、\(2\) 号同学当选的指示量分别为 \(a_{i1}\)、\(a_{i2}\). 由于它们只取 \(0\) 或 \(1\)，两者的乘积在且仅在二者都为 \(1\) 时等于 \(1\)，所以同时同意的人数为 \(\sum_{i=1}^k a_{i1}a_{i2}=a_{11}a_{12}+a_{21}a_{22}+\cdots+a_{k1}a_{k2}\)，选 C.
\end{solution}

\section{填空题}
\begin{problem}
已知某球体的体积与其表面积的数值相等，则此球体的半径为\fillinblank{}.
\end{problem}

\begin{answer}
\(3\)
\end{answer}

\begin{solution}
设球半径为 \(r>0\). 球体积为 \(V=4\pi r^3/3\)，表面积为 \(S=4\pi r^2\). 由数值相等得 \(4\pi r^3/3=4\pi r^2\)，约去 \(4\pi r^2\) 后得 \(r/3=1\)，所以 \(r=3\).
\end{solution}

\begin{problem}
函数 \(f(x) = \lg(1 + x^2)\)，\(g(x) = 2 - |x|\)，\(h(x) = \tan 2x\) 中，\fillinblank{}是偶函数.
\end{problem}

\begin{answer}
\(f(x)\)，\(g(x)\)
\end{answer}

\begin{solution}
有 \(f(-x)=\lg(1+(-x)^2)=\lg(1+x^2)=f(x)\)，所以 \(f\) 是偶函数；又 \(g(-x)=2-|-x|=2-|x|=g(x)\)，所以 \(g\) 是偶函数. 对 \(h\)，在其定义域上 \(h(-x)=\tan(-2x)=-\tan2x=-h(x)\)，它是奇函数而不是偶函数，故应填 \(f(x)\)，\(g(x)\).
\end{solution}

\begin{problem}
以双曲线 \(\frac{x^{2}}{16} - \frac{y^{2}}{9} = 1\) 右顶点为顶点，左焦点为焦点的抛物线的方程是\fillinblank{}.
\end{problem}

\begin{answer}
\(y^2=-36(x-4)\)
\end{answer}

\begin{solution}
双曲线的 \(a=4\)，\(b=3\)，故 \(c=\sqrt{a^2+b^2}=5\). 其右顶点为 \((4,0)\)，左焦点为 \((-5,0)\). 以 \((4,0)\) 为顶点、焦点为 \((-5,0)\) 的抛物线向左开口，设方程为 \(y^2=4p(x-4)\). 顶点到焦点的有向距离为 \(p=-9\)，所以方程为 \(y^2=-36(x-4)\).
\end{solution}

\begin{problem}
将长度为 \(1\) 的铁丝分成两段，分别围成一个正方形和一个圆形，要使正方形与圆的面积之和最小，正方形的周长应为\fillinblank{}.
\end{problem}

\begin{answer}
\(\frac{4}{\pi+4}\)
\end{answer}

\begin{solution}
设正方形周长为 \(x\)，则圆的周长为 \(1-x\)，其中 \(0<x<1\). 正方形边长为 \(x/4\)，圆半径为 \((1-x)/(2\pi)\)，所以面积和为
\(F(x)=\dfrac{x^2}{16}+\dfrac{(1-x)^2}{4\pi}\). 求导得
\(F'(x)=\dfrac{x}{8}-\dfrac{1-x}{2\pi}\). 令 \(F'(x)=0\)，得 \(\pi x=4(1-x)\)，即 \(x=\dfrac{4}{\pi+4}\)，且此值属于 \((0,1)\). 又
\(F''(x)=\dfrac18+\dfrac{1}{2\pi}>0\)，故 \(F\) 在该点取得最小值，正方形周长为 \(\dfrac{4}{\pi+4}\).
\end{solution}

\section{解答题}
\begin{problem}
已知函数 \(f(x) = \cos^4 x - 2 \sin x \cos x - \sin^4 x\).

\begin{enumerate}
\item 求 \(f(x)\) 的最小正周期；
\item 求 \(f(x)\) 的最大值，最小值.
\end{enumerate}
\end{problem}

\begin{solution}
\begin{enumerate}
\item 由平方差公式和二倍角公式，
\(f(x)=\bigl(\cos^4x-\sin^4x\bigr)-2\sin x\cos x=(\cos^2x-\sin^2x)-\sin2x=\cos2x-\sin2x=\sqrt2\cos\left(2x+\dfrac\pi4\right)\). 若 \(T>0\) 是其周期，则 \(2T\) 必为余弦函数的周期，即 \(2T=2k\pi\)；因此最小的正周期为 \(T=\pi\).
\item 因为 \(-1\leqslant\cos(2x+\pi/4)\leqslant1\)，所以 \(-\sqrt2\leqslant f(x)\leqslant\sqrt2\)，且余弦可以取到 \(1\) 与 \(-1\). 故 \(f(x)\) 的最大值为 \(\sqrt2\)，最小值为 \(-\sqrt2\).
\end{enumerate}
\end{solution}

\begin{problem}
已知数列 \(\{a_n\}\) 是等差数列，且 \(a_1 = 2\)，\(a_1 + a_2 + a_3 = 12\).

\begin{enumerate}
\item 求数列 \(\{a_n\}\) 的通项公式；
\item 令 \(b_n = a_n \cdot 3^n\)，求数列 \(\{b_n\}\) 前 \(n\) 项和的公式.
\end{enumerate}
\end{problem}

\begin{solution}
\begin{enumerate}
\item 设公差为 \(d\). 由等差数列三项和公式，\(a_1+a_2+a_3=3a_1+3d=12\)，代入 \(a_1=2\) 得 \(6+3d=12\)，所以 \(d=2\)，从而 \(a_n=a_1+(n-1)d=2n\).
\item 由 \(b_k=a_k3^k=2k3^k\)，记前 \(n\) 项和为 \(S_n\)，则 \(S_n=2\cdot3+4\cdot3^2+\cdots+2n\cdot3^n\). 两边乘以 \(3\)，得 \(3S_n=2\cdot3^2+4\cdot3^3+\cdots+2(n-1)\cdot3^n+2n\cdot3^{n+1}\). 两式相减，得到

\[
2S_n=2n\cdot3^{n+1}-2(3+3^2+\cdots+3^n)=2n\cdot3^{n+1}-3(3^n-1)=3\bigl((2n-1)3^n+1\bigr)
\]

所以 \(S_n=\dfrac{3}{2}\bigl((2n-1)3^n+1\bigr)\).
\end{enumerate}
\end{solution}

\begin{problem}
如图，正三棱柱 \(ABC - A_1B_1C_1\) 中，\(D\) 是 \(BC\) 的中点，\(AB = a\).

\begin{enumerate}
\item 求证：直线 \(A_1D \perp B_1C_1\)；
\item 求点 \(D\) 到平面 \(ACC_1\) 的距离；
\item 判断 \(A_1B\) 与平面 \(ADC_1\) 的位置关系，并证明你的结论.
\end{enumerate}

\bitmapfigure[max width=0.42\linewidth]{img/2003/beijing_liberal_autumn/q17_fig1.png}
\end{problem}

\begin{solution}
\begin{enumerate}
\item 设 \(E\) 为 \(B_1C_1\) 的中点. 在正三角形 \(A_1B_1C_1\) 中，\(A_1E\perp B_1C_1\). 又因 \(D\)、\(E\) 分别为 \(BC\)、\(B_1C_1\) 的中点，且棱柱为直棱柱，所以 \(DE\parallel BB_1\perp B_1C_1\). 直线 \(A_1E\)、\(DE\) 相交且都在平面 \(A_1DE\) 内，故 \(B_1C_1\perp\text{平面 }A_1DE\)，从而 \(A_1D\perp B_1C_1\).
\item 设 \(H\) 为 \(D\) 到 \(AC\) 的垂足. 因为 \(AC\subset\text{平面 }ACC_1\)，且 \(CC_1\perp\text{平面 }ABC\)，所以 \(AC\) 与 \(CC_1\) 是平面 \(ACC_1\) 内的两条相交直线. 又 \(DH\perp AC\)，\(DH\subset\text{平面 }ABC\)，故 \(DH\perp CC_1\)，从而 \(DH\perp\text{平面 }ACC_1\). 在等边三角形 \(ABC\) 中，\(D\) 是 \(BC\) 中点，三角形 \(ADC\) 面积为三角形 \(ABC\) 面积的一半，故 \(DH=\dfrac{2S_{\triangle ADC}}{AC}=\dfrac{\sqrt3}{4}a\). 这就是所求距离.
\item 在平行四边形 \(ACC_1A_1\) 中，设对角线 \(A_1C\) 与 \(AC_1\) 交于 \(F\)，则 \(F\) 是 \(A_1C\) 的中点. 在三角形 \(A_1CB\) 中，\(F\)，\(D\) 分别是 \(A_1C\)，\(CB\) 的中点，所以 \(DF\parallel A_1B\). 又 \(DF\subset\text{平面 }ADC_1\). 下面说明 \(A_1B\) 不在该平面内：若 \(A_1\) 在平面 \(ADC_1\) 内，则 \(AA_1\) 在该平面内；由于 \(AA_1\parallel CC_1\) 且 \(C_1\) 在该平面内，点 \(C\) 也在该平面内，从而该平面就是底面 \(ABC\)，这与 \(C_1\) 不在底面 \(ABC\) 内矛盾. 因此 \(A_1B\parallel\text{平面 }ADC_1\).
\end{enumerate}
\end{solution}

\begin{problem}
如图，\(A_1\)，\(A_2\) 为椭圆的两个顶点，\(F_1\)，\(F_2\) 为椭圆的两个焦点.

\begin{enumerate}
\item 写出椭圆的方程及准线方程；
\item 过线段 \(OA\) 上异于 \(O\)，\(A\) 的任一点 \(K\) 作 \(OA\) 的垂线，交椭圆于 \(P\)，\(P_1\) 两点，直线 \(A_1P\) 与 \(AP_1\) 交于点 \(M\). 求证：点 \(M\) 在双曲线 \(\frac{x^{2}}{25} - \frac{y^{2}}{9} = 1\) 上.
\end{enumerate}

\bitmapfigure[max width=0.62\linewidth]{img/2003/beijing_liberal_autumn/q18_fig1.png}
\end{problem}

\begin{solution}
\begin{enumerate}
\item 由图中顶点 \(A_1=(-5,0)\)，\(A=(5,0)\) 得长半轴 \(a=5\)，由焦点 \(F_1=(-4,0)\)，\(F_2=(4,0)\) 得焦半距 \(c=4\)，所以 \(b^2=a^2-c^2=25-16=9\). 椭圆方程为 \(x^2/25+y^2/9=1\)，离心率为 \(4/5\)，故准线方程为 \(x=\pm a/e=\pm25/4\).
\item 设 \(K=(x_0,0)\)，其中 \(0<x_0<5\)，直线 \(x=x_0\) 与椭圆交于 \(P=(x_0,y_0)\)、\(P_1=(x_0,-y_0)\). 直线 \(A_1P\)、\(AP_1\) 的方程分别为 \(y=y_0(x+5)/(x_0+5)\) 与 \(y=-y_0(x-5)/(x_0-5)\). 设交点 \(M=(x,y)\)，联立两式并约去 \(y_0\ne0\)，得 \(x=25/x_0\)，再代回得 \(y=5y_0/x_0\). 由 \(P\) 在椭圆上，\(x_0^2/25+y_0^2/9=1\)，于是 \(x^2/25-y^2/9=25/x_0^2-25y_0^2/(9x_0^2)=25/x_0^2(1-y_0^2/9)=25/x_0^2\cdot x_0^2/25=1\). 因此 \(M\) 在双曲线 \(x^2/25-y^2/9=1\) 上.
\end{enumerate}
\end{solution}

\begin{problem}
有三个新兴城镇，分别位于 \(A\)，\(B\)，\(C\) 三点处，且 \(AB = AC = 13 \mathrm{~km}\)，\(BC = 10 \mathrm{~km}\). 今计划合建一个中心医院，为同时方便三镇，准备建在 \(BC\) 的垂直平分线上的 \(P\) 点处.（建立坐标系如图）

\begin{enumerate}
\item 若希望点 \(P\) 到三镇距离的平方和为最小，点 \(P\) 应位于何处？
\item 若希望点 \(P\) 到三镇的最远距离为最小，点 \(P\) 应位于何处？
\end{enumerate}

\bitmapfigure[max width=0.44\linewidth]{img/2003/beijing_liberal_autumn/q19_fig1.png}
\end{problem}

\begin{solution}
\begin{enumerate}
\item 由 \(AB=AC=13\)、\(B=(-5,0)\)、\(C=(5,0)\)，得 \(A=(0,12)\). 设 \(P=(0,y)\)，则三镇距离的平方和为 \(PA^2+PB^2+PC^2=(y-12)^2+(y^2+25)+(y^2+25)=3(y-4)^2+146\). 该式在 \(y=4\) 时取得最小值，所以 \(P=(0,4)\).
\item 设 \(g(y)\) 为点 \(P=(0,y)\) 到三镇的最远距离，则 \(g(y)=\max\{\sqrt{y^2+25},|y-12|\}\). 当 \(y\leqslant12\) 时，比较两者平方：\(y^2+25-(12-y)^2=24y-119\). 因而 \(y<119/24\) 时 \(g(y)=12-y\)，它随 \(y\) 增大而减小；\(119/24\leqslant y\leqslant12\) 时 \(g(y)=\sqrt{y^2+25}\)，它随 \(y\) 增大而增大. 当 \(y\geqslant12\) 时，\(y^2+25-(y-12)^2=24y-119>0\)，仍有 \(g(y)=\sqrt{y^2+25}\)，并随 \(y\) 增大而增大. 所以最小值在两段交界处取得，即 \(P=\left(0,\dfrac{119}{24}\right)\).
\end{enumerate}
\end{solution}

\begin{problem}
设 \(y = f(x)\) 是定义在区间 \([-1,1]\) 上的函数，且满足条件：

\begin{circlelist}
\item \(f(-1) = f(1) = 0\)；
\item 对任意的 \(u\)，\(v \in [-1,1]\)，都有 \(\left|f(u) - f(v)\right| \leqslant |u - v|\).
\end{circlelist}

\begin{enumerate}
\item 证明：对任意的 \(x \in [-1,1]\)，\(x - 1 \leqslant f(x) \leqslant 1 - x\)；
\item 判断函数 \(g(x) = \begin{cases}
1 + x,& x \in [-1,0),\\
1 - x,& x \in [0,1]
\end{cases}\) 是否满足题设条件；
\item 在区间 \([-1,1]\) 上是否存在满足题设条件的函数 \(y = f(x)\)，且使得对任意的 \(u\)，\(v \in [-1,1]\)，都有 \(\left|f(u) - f(v)\right| = u - v\)，若存在，请举一例；若不存在，请说明理由.
\end{enumerate}
\end{problem}

\begin{solution}
\begin{enumerate}
\item 取题设不等式中的 \(u=x\)，\(v=1\)，由 \(f(1)=0\) 得 \(|f(x)|=|f(x)-f(1)|\leqslant|x-1|=1-x\). 因而 \(-(1-x)\leqslant f(x)\leqslant1-x\)，即 \(x-1\leqslant f(x)\leqslant1-x\).
\item 先有 \(g(-1)=0=g(1)\). 若 \(u\)，\(v\in[0,1]\)，则 \(|g(u)-g(v)|=|(1-u)-(1-v)|=|u-v|\)；若 \(u\)，\(v\in[-1,0)\)，则 \(|g(u)-g(v)|=|(1+u)-(1+v)|=|u-v|\). 不妨设 \(u\in[-1,0)\)，\(v\in[0,1]\)，则 \(|g(u)-g(v)|=|u+v|\)，而 \(u\leqslant0\leqslant v\) 时 \((v-u)^2-(u+v)^2=-4uv\geqslant0\)，故 \(|u+v|\leqslant v-u=|u-v|\). 因此 \(g\) 满足题设条件.
\item 不存在. 若存在这样的函数，取 \(u=-1\)，\(v=1\)，则题设等式给出 \(|f(-1)-f(1)|=-2\). 左边等于 \(0\)，右边等于 \(-2\)，矛盾，所以不存在.
\end{enumerate}
\end{solution}
